% !TEX root = ../main.tex
% =====================================================================
%  PREAMBULE – všechny balíčky, písmo a globální nastavení sazby.
%  Tento soubor se načítá z main.tex; běžně do něj nemusíte zasahovat.
% =====================================================================
\documentclass[12pt,a4paper]{article}

% ---------- geometrie stránky ----------
\usepackage[a4paper,
            left=2.5cm, right=2.4cm,
            top=2.0cm, bottom=2.5cm,
            footskip=16pt]{geometry}

% ---------- písmo a jazyk ----------
\usepackage{fontspec}
\usepackage{amsmath}
\usepackage{unicode-math}
\IfFontExistsTF{Times New Roman}
  {\setmainfont{Times New Roman}}
  {\setmainfont{TeX Gyre Termes}}
\IfFontExistsTF{Cambria Math}
  {\setmathfont{Cambria Math}}
  {\IfFontExistsTF{TeX Gyre Termes Math}
    {\setmathfont{TeX Gyre Termes Math}}
    {\setmathfont{Latin Modern Math}}}
\usepackage{polyglossia}
\setdefaultlanguage{czech}

% ---------- základní sazba ----------
\linespread{0.95}
\setlength{\parindent}{0.6cm}
\setlength{\parskip}{6pt}
\setlength{\emergencystretch}{3em}
\tolerance=1200
\hyphenpenalty=10000
\exhyphenpenalty=10000
\setlength{\abovedisplayskip}{8pt}
\setlength{\belowdisplayskip}{8pt}
\setlength{\abovedisplayshortskip}{4pt}
\setlength{\belowdisplayshortskip}{4pt}

\usepackage{xcolor}
\usepackage{graphicx}
\graphicspath{{obrazky/}{./}}
\usepackage{float}
\usepackage{needspace}
\usepackage[font=small,labelfont=bf,justification=centering,skip=6pt]{caption}
\usepackage{array}
\usepackage{multirow}
\usepackage{enumitem}
\usepackage{titlesec}
\usepackage{fancyhdr}
\usepackage{listings}
\usepackage[hidelinks]{hyperref}

% ---------- nadpisy ----------
\titleformat{\section}[hang]
  {\normalfont\fontsize{14}{16.8}\selectfont\bfseries}
  {\makebox[0.63cm][l]{\thesection.}}{0pt}{}
\titlespacing*{\section}{0pt}{14pt}{2pt}
\titleformat{\subsection}[hang]
  {\normalfont\fontsize{13}{15.6}\selectfont\bfseries}
  {}{0pt}{}
\titlespacing*{\subsection}{0pt}{12pt}{2pt}
\setcounter{secnumdepth}{1}

% ---------- seznamy ----------
\setlist[itemize]{leftmargin=1.27cm, labelindent=0.63cm,
  labelwidth=0.63cm, labelsep=0pt, align=left,
  topsep=2pt, itemsep=2pt, parsep=0pt, partopsep=0pt}
\setlist[enumerate]{leftmargin=1.27cm, labelindent=0.63cm,
  labelwidth=0.63cm, labelsep=0pt, align=left,
  topsep=2pt, itemsep=2pt, parsep=0pt, partopsep=0pt}

% ---------- zápatí ----------
\pagestyle{fancy}
\fancyhf{}
\fancyfoot[C]{- \thepage\ -}
\renewcommand{\headrulewidth}{0pt}
\renewcommand{\footrulewidth}{0pt}
\fancypagestyle{plain}{%
  \fancyhf{}%
  \fancyfoot[C]{- \thepage\ -}%
  \renewcommand{\headrulewidth}{0pt}%
  \renewcommand{\footrulewidth}{0pt}%
}

% ---------- české názvy ----------
\captionsetup[figure]{name=Obr.}
\captionsetup[table]{name=Tab.}
\renewcommand{\figurename}{Obr.}
\renewcommand{\tablename}{Tab.}
\renewcommand{\refname}{Použitá literatura}
\renewcommand{\contentsname}{Obsah}
\renewcommand{\listfigurename}{Seznam obrázků}
\renewcommand{\listtablename}{Seznam tabulek}
\setlength{\belowcaptionskip}{2pt}

% Některé názvy nastavuje polyglossia až na začátku dokumentu.
\AtBeginDocument{%
  \renewcommand{\refname}{Použitá literatura}%
  \renewcommand{\lstlistingname}{Výpis}%
}

% ---------- výpisy kódu ----------
\lstset{
  basicstyle=\ttfamily\small,
  breaklines=true,
  frame=single,
  rulecolor=\color{gray!50},
  backgroundcolor=\color{gray!5},
  keywordstyle=\bfseries,
  commentstyle=\itshape,
  stringstyle=\ttfamily,
  showstringspaces=false,
  tabsize=2,
  columns=fullflexible,
  keepspaces=true
}
